
% =========================================================
% Supplementary Chapter: Simple Harmonic Motion
% POSN 1 Physics Preparation Notes
% Language: Thai
% Compiler: XeLaTeX
% =========================================================

\chapter{บทเสริม: การเคลื่อนที่แบบฮาร์มอนิกอย่างง่าย}

\begin{center}
    {\Large \textbf{Simple Harmonic Motion}}\\[4pt]
    {\normalsize เอกสารเตรียมสอบคัดเลือก สอวน. ฟิสิกส์ ค่าย 1}\\[6pt]
    {\normalsize จัดทำโดย \textbf{Tames Thanakrit Damduan}}\\[2pt]
    {\footnotesize \textcopyright\ 2026 Tames Thanakrit Damduan. All rights reserved.}\\[8pt]
\end{center}


\section*{เป้าหมายของบทเรียน}

หลังจากเรียนบทนี้ นักเรียนควรทำสิ่งต่อไปนี้ได้

\begin{enumerate}
    \item เข้าใจความหมายพื้นฐานของการเคลื่อนที่แบบฮาร์มอนิกอย่างง่าย
    \item ใช้ความสัมพันธ์ระหว่างคาบ ความถี่ และความถี่เชิงมุมได้
    \item ใช้สูตรคาบของมวลติดสปริงได้
    \item ใช้สูตรคาบของลูกตุ้มอย่างง่ายเมื่อมุมแกว่งมีค่าน้อยได้
    \item เปรียบเทียบคาบก่อนและหลังจากตัวแปรบางตัวเปลี่ยนแปลงได้
    \item ประยุกต์สูตรคาบของลูกตุ้มร่วมกับการขยายตัวเชิงความร้อนได้
\end{enumerate}

\begin{tcolorbox}[title=แนวคิดสำคัญของบทนี้]
หัวใจของการเคลื่อนที่แบบฮาร์มอนิกอย่างง่าย คือ เมื่อวัตถุถูกเบนออกจากตำแหน่งสมดุล จะมีแรงหรือความเร่งพยายามดึงวัตถุกลับเข้าหาตำแหน่งสมดุล
\end{tcolorbox}

% =========================================================
\section{การเคลื่อนที่แบบฮาร์มอนิกอย่างง่าย}

\subsection{ตำแหน่งสมดุล}

ตำแหน่งสมดุล คือ ตำแหน่งที่แรงลัพธ์บนวัตถุเป็นศูนย์

\[
\sum F=0
\]

ถ้าวัตถุถูกดึงออกจากตำแหน่งสมดุล แล้วมีแรงพยายามดึงวัตถุกลับมา แรงนี้เรียกว่า แรงคืนรูป

\[
\text{restoring force}
\]

\subsection{เงื่อนไขของ SHM}

ถ้าเลือกตำแหน่งสมดุลเป็น

\[
x=0
\]

การเคลื่อนที่แบบฮาร์มอนิกอย่างง่าย หรือ SHM จะมีความเร่งเป็น

\[
a=-\omega^2x
\]

โดย

\[
x=\text{การกระจัดจากตำแหน่งสมดุล}
\]

\[
\omega=\text{ความถี่เชิงมุม}
\]

เครื่องหมายลบหมายความว่า ความเร่งชี้กลับเข้าหาตำแหน่งสมดุลเสมอ

ถ้า

\[
x>0
\]

จะได้

\[
a<0
\]

ถ้า

\[
x<0
\]

จะได้

\[
a>0
\]

\begin{tcolorbox}[title=เงื่อนไขของ SHM]
\[
a=-\omega^2x
\]

ความเร่งแปรผันตรงกับการกระจัด แต่มีทิศตรงข้ามกับการกระจัด
\end{tcolorbox}

% =========================================================
\section{คาบ ความถี่ และความถี่เชิงมุม}

\subsection{คาบ}

คาบ คือ เวลาที่วัตถุใช้เคลื่อนที่ครบหนึ่งรอบของการสั่น

ใช้สัญลักษณ์

\[
T
\]

หน่วยคือ

\[
\si{s}
\]

\subsection{ความถี่}

ความถี่ คือ จำนวนรอบของการสั่นต่อหนึ่งหน่วยเวลา

ใช้สัญลักษณ์

\[
f
\]

หน่วยคือ

\[
\si{Hz}
\]

ความถี่และคาบมีความสัมพันธ์กันเป็น

\[
f=\frac{1}{T}
\]

\[
T=\frac{1}{f}
\]

\subsection{ความถี่เชิงมุม}

หนึ่งรอบของการสั่นเทียบเท่ากับมุม

\[
2\pi
\]

เรเดียน ดังนั้นความถี่เชิงมุมคือ

\[
\omega=2\pi f
\]

เพราะ

\[
f=\frac{1}{T}
\]

จึงได้

\[
\omega=\frac{2\pi}{T}
\]

และ

\[
T=\frac{2\pi}{\omega}
\]

\begin{tcolorbox}[title=ความสัมพันธ์พื้นฐาน]
\[
f=\frac{1}{T}
\]

\[
\omega=2\pi f
\]

\[
\omega=\frac{2\pi}{T}
\]

\[
T=\frac{2\pi}{\omega}
\]
\end{tcolorbox}

% =========================================================
\section{มวลติดสปริง}

\subsection{กฎของฮุก}

สำหรับสปริงอุดมคติ ถ้ามวลถูกดึงหรืออัดจากตำแหน่งสมดุลเป็นระยะ $x$ แรงของสปริงคือ

\[
F_s=-kx
\]

โดย

\[
k=\text{ค่าคงตัวสปริง}
\]

หน่วยของ $k$ คือ

\[
\si{N/m}
\]

เครื่องหมายลบหมายถึงแรงของสปริงชี้กลับเข้าหาตำแหน่งสมดุล

\begin{tcolorbox}[title=กฎของฮุก]
\[
F_s=-kx
\]
\end{tcolorbox}

\subsection{คาบของมวลติดสปริง}

ให้มวล $m$ ติดกับสปริงและเคลื่อนที่บนพื้นลื่น

ใช้กฎข้อที่สองของนิวตัน

\[
\sum F=ma
\]

แรงลัพธ์คือแรงของสปริง

\[
ma=-kx
\]

จัดรูป

\[
a=-\frac{k}{m}x
\]

เทียบกับเงื่อนไขของ SHM

\[
a=-\omega^2x
\]

จะได้

\[
\omega^2=\frac{k}{m}
\]

ดังนั้น

\[
\omega=\sqrt{\frac{k}{m}}
\]

และเพราะ

\[
T=\frac{2\pi}{\omega}
\]

จึงได้

\[
T=2\pi\sqrt{\frac{m}{k}}
\]

\begin{tcolorbox}[title=มวลติดสปริง]
\[
\omega=\sqrt{\frac{k}{m}}
\]

\[
T=2\pi\sqrt{\frac{m}{k}}
\]
\end{tcolorbox}

\subsection{การเปรียบเทียบคาบของมวลติดสปริง}

จาก

\[
T=2\pi\sqrt{\frac{m}{k}}
\]

จะเห็นว่า

\[
T\propto \sqrt{m}
\]

และ

\[
T\propto \frac{1}{\sqrt{k}}
\]

ดังนั้น

\[
\frac{T_2}{T_1}
=
\sqrt{
\frac{m_2}{m_1}
\frac{k_1}{k_2}
}
\]

\begin{examplebox}{ตัวอย่าง: เปรียบเทียบคาบของสปริง}
มวลติดสปริงมีคาบ $T$ ถ้าเปลี่ยนมวลจาก $m$ เป็น $4m$ โดยใช้สปริงเดิม จงหาคาบใหม่

เริ่มจาก

\[
T\propto \sqrt{m}
\]

ดังนั้น

\[
\frac{T'}{T}
=
\sqrt{\frac{4m}{m}}
\]

\[
\frac{T'}{T}=2
\]

จึงได้

\[
\boxed{T'=2T}
\]
\end{examplebox}

% =========================================================
\section{ลูกตุ้มอย่างง่าย}

\subsection{องค์ประกอบของลูกตุ้มอย่างง่าย}

ลูกตุ้มอย่างง่ายประกอบด้วย

\begin{enumerate}
    \item มวลขนาดเล็ก
    \item เชือกเบาที่ไม่ยืด
    \item จุดแขวนที่อยู่นิ่ง
\end{enumerate}

ให้ความยาวเชือกเป็น

\[
\ell
\]

เมื่อลูกตุ้มถูกเบนจากแนวดิ่งเป็นมุม $\theta$ น้ำหนักจะมีองค์ประกอบตามแนวสัมผัสเป็น

\[
F_t=-mg\sin\theta
\]

เครื่องหมายลบหมายถึงแรงนี้พยายามดึงลูกตุ้มกลับเข้าหาตำแหน่งสมดุล

\subsection{เงื่อนไขมุมเล็ก}

ถ้ามุม $\theta$ มีค่าน้อย และวัดในหน่วยเรเดียน สามารถใช้ค่าประมาณได้ว่า

\[
\sin\theta\approx\theta
\]

ระยะตามแนวส่วนโค้งคือ

\[
s=\ell\theta
\]

ดังนั้น

\[
\theta=\frac{s}{\ell}
\]

เริ่มจากแรงตามแนวสัมผัส

\[
F_t=-mg\sin\theta
\]

สำหรับมุมเล็ก

\[
F_t\approx -mg\theta
\]

แทน

\[
\theta=\frac{s}{\ell}
\]

จะได้

\[
F_t=-\frac{mg}{\ell}s
\]

ใช้กฎข้อที่สองของนิวตัน

\[
ma_t=-\frac{mg}{\ell}s
\]

ตัด $m$

\[
a_t=-\frac{g}{\ell}s
\]

เทียบกับเงื่อนไขของ SHM

\[
a=-\omega^2x
\]

จะได้

\[
\omega^2=\frac{g}{\ell}
\]

ดังนั้น

\[
\omega=\sqrt{\frac{g}{\ell}}
\]

และ

\[
T=\frac{2\pi}{\omega}
\]

จึงได้

\[
T=2\pi\sqrt{\frac{\ell}{g}}
\]

\begin{tcolorbox}[title=คาบของลูกตุ้มอย่างง่าย]
สำหรับมุมแกว่งที่มีค่าน้อย

\[
T=2\pi\sqrt{\frac{\ell}{g}}
\]
\end{tcolorbox}

\subsection{สิ่งที่คาบของลูกตุ้มขึ้นกับ}

จาก

\[
T=2\pi\sqrt{\frac{\ell}{g}}
\]

จะเห็นว่า

\[
T\propto \sqrt{\ell}
\]

และ

\[
T\propto \frac{1}{\sqrt{g}}
\]

แต่คาบไม่ขึ้นกับมวลของลูกตุ้ม

\[
T\not\propto m
\]

สำหรับมุมแกว่งขนาดเล็ก คาบไม่ขึ้นกับขนาดของมุมเริ่มต้นโดยประมาณ

\begin{tcolorbox}[title=ข้อควรจำ]
คาบของลูกตุ้มอย่างง่ายขึ้นกับ

\[
\ell
\]

และ

\[
g
\]

แต่ไม่ขึ้นกับมวล
\end{tcolorbox}

\subsection{เปรียบเทียบคาบของลูกตุ้ม}

ถ้าความยาวหรือความเร่งโน้มถ่วงเปลี่ยนไป ใช้อัตราส่วน

\[
\frac{T_2}{T_1}
=
\sqrt{
\frac{\ell_2}{\ell_1}
\frac{g_1}{g_2}
}
\]

\begin{examplebox}{ตัวอย่าง: เปลี่ยนความยาวเชือก}
ลูกตุ้มอย่างง่ายมีคาบ $T$ ถ้าเพิ่มความยาวเชือกจาก $\ell$ เป็น $4\ell$ โดยค่า $g$ คงเดิม จงหาคาบใหม่

เริ่มจาก

\[
T\propto\sqrt{\ell}
\]

ดังนั้น

\[
\frac{T'}{T}
=
\sqrt{\frac{4\ell}{\ell}}
\]

\[
\frac{T'}{T}=2
\]

จึงได้

\[
\boxed{T'=2T}
\]
\end{examplebox}

% =========================================================
\section{ลูกตุ้มกับการขยายตัวเชิงความร้อน}

หัวข้อนี้เชื่อมกับเรื่องการขยายตัวเชิงความร้อน ซึ่งพบในขอบเขตเนื้อหาสอบ

ถ้าลวดยาวเดิม

\[
\ell
\]

เมื่ออุณหภูมิเพิ่มขึ้น

\[
\Delta T
\]

และสัมประสิทธิ์การขยายตัวเชิงเส้นเป็น

\[
\alpha
\]

ความยาวใหม่คือ

\[
\ell'=\ell(1+\alpha\Delta T)
\]

สำหรับลูกตุ้มอย่างง่าย

\[
T=2\pi\sqrt{\frac{\ell}{g}}
\]

ดังนั้น

\[
\frac{T'}{T}
=
\sqrt{\frac{\ell'}{\ell}}
\]

แทน

\[
\ell'=\ell(1+\alpha\Delta T)
\]

จะได้

\[
\frac{T'}{T}
=
\sqrt{1+\alpha\Delta T}
\]

\begin{tcolorbox}[title=ลูกตุ้มกับการขยายตัวเชิงความร้อน]
\[
\ell'=\ell(1+\alpha\Delta T)
\]

\[
\frac{T'}{T}
=
\sqrt{1+\alpha\Delta T}
\]
\end{tcolorbox}

% =========================================================
\section{ชุดโจทย์รวมท้ายบท}

\subsection*{ระดับพื้นฐาน}

\begin{enumerate}
    \item วัตถุหนึ่งสั่นด้วยคาบ $0.50\,\si{s}$ จงหาความถี่
    \item วัตถุหนึ่งมีความถี่ $4\,\si{Hz}$ จงหาคาบ
    \item วัตถุหนึ่งมีคาบ $2\pi\,\si{s}$ จงหาความถี่เชิงมุม
    \item มวล $1\,\si{kg}$ ติดกับสปริงที่มี $k=100\,\si{N/m}$ จงหาคาบ
\end{enumerate}

\subsection*{ระดับกลาง}

\begin{enumerate}
    \setcounter{enumi}{4}
    \item มวลติดสปริงมีคาบ $T$ ถ้าเพิ่มมวลเป็น $9$ เท่าของเดิม คาบใหม่เป็นกี่เท่าของเดิม
    \item มวลติดสปริงมีคาบ $T$ ถ้าเปลี่ยนค่าคงตัวสปริงเป็น $4k$ คาบใหม่เป็นกี่เท่าของเดิม
    \item ลูกตุ้มอย่างง่ายมีคาบ $T$ ถ้าเพิ่มความยาวเชือกเป็น $9$ เท่าของเดิม คาบใหม่เป็นกี่เท่าของเดิม
    \item ลูกตุ้มอย่างง่ายมีคาบ $T$ บนโลก ถ้านำไปอยู่ในบริเวณที่มี $g'=g/4$ คาบใหม่เป็นกี่เท่าของเดิม
\end{enumerate}

\subsection*{ระดับประยุกต์}

\begin{enumerate}
    \setcounter{enumi}{8}
    \item ลูกตุ้มอย่างง่ายยาว $\ell$ มีคาบ $T$ ถ้าอุณหภูมิเพิ่มขึ้น $\Delta T$ และสัมประสิทธิ์การขยายตัวเชิงเส้นของลวดเป็น $\alpha$ จงหาอัตราส่วน $T'/T$
    \item ลูกตุ้มสองลูกมีความยาวเท่ากัน ถ้าลูกตุ้มทั้งสองผ่านตำแหน่งต่ำสุดพร้อมกัน แล้วแยกออกจากกัน จงหาเวลาที่ทั้งสองกลับมาถึงตำแหน่งต่ำสุดพร้อมกันอีกครั้ง
\end{enumerate}

% =========================================================
\section{เฉลยย่อท้ายบท}

\begin{enumerate}
    \item
    \[
    f=\frac{1}{T}
    \]
    \[
    f=\frac{1}{0.50}=2\,\si{Hz}
    \]

    \item
    \[
    T=\frac{1}{f}
    \]
    \[
    T=\frac{1}{4}=0.25\,\si{s}
    \]

    \item
    \[
    \omega=\frac{2\pi}{T}
    \]
    \[
    \omega=\frac{2\pi}{2\pi}=1\,\si{rad/s}
    \]

    \item
    \[
    T=2\pi\sqrt{\frac{m}{k}}
    \]
    \[
    T=2\pi\sqrt{\frac{1}{100}}
    \]
    \[
    T=\frac{\pi}{5}\,\si{s}
    \]

    \item
    \[
    T\propto\sqrt{m}
    \]
    \[
    \frac{T'}{T}=\sqrt{9}=3
    \]
    \[
    \boxed{T'=3T}
    \]

    \item
    \[
    T\propto\frac{1}{\sqrt{k}}
    \]
    \[
    \frac{T'}{T}=\sqrt{\frac{k}{4k}}=\frac{1}{2}
    \]
    \[
    \boxed{T'=\frac{T}{2}}
    \]

    \item
    \[
    T\propto\sqrt{\ell}
    \]
    \[
    \frac{T'}{T}=\sqrt{9}=3
    \]
    \[
    \boxed{T'=3T}
    \]

    \item
    \[
    T\propto\frac{1}{\sqrt{g}}
    \]
    \[
    \frac{T'}{T}
    =
    \sqrt{\frac{g}{g/4}}
    \]
    \[
    \frac{T'}{T}=2
    \]
    \[
    \boxed{T'=2T}
    \]

    \item
    \[
    \ell'=\ell(1+\alpha\Delta T)
    \]
    \[
    \frac{T'}{T}
    =
    \sqrt{\frac{\ell'}{\ell}}
    \]
    \[
    \boxed{
    \frac{T'}{T}
    =
    \sqrt{1+\alpha\Delta T}
    }
    \]

    \item
    ลูกตุ้มแต่ละลูกกลับมาที่ตำแหน่งต่ำสุดอีกครั้งหลังเวลาครึ่งคาบ
    \[
    \boxed{\Delta t=\frac{T}{2}}
    \]
\end{enumerate}

% =========================================================
\section{สรุปสูตรสำคัญ}

\begin{tcolorbox}[title=สรุปบทเสริม: การเคลื่อนที่แบบฮาร์มอนิกอย่างง่าย]
เงื่อนไขของ SHM

\[
a=-\omega^2x
\]

คาบ ความถี่ และความถี่เชิงมุม

\[
f=\frac{1}{T}
\]

\[
\omega=2\pi f
\]

\[
\omega=\frac{2\pi}{T}
\]

มวลติดสปริง

\[
F_s=-kx
\]

\[
\omega=\sqrt{\frac{k}{m}}
\]

\[
T=2\pi\sqrt{\frac{m}{k}}
\]

ลูกตุ้มอย่างง่ายสำหรับมุมเล็ก

\[
\sin\theta\approx\theta
\]

\[
\omega=\sqrt{\frac{g}{\ell}}
\]

\[
T=2\pi\sqrt{\frac{\ell}{g}}
\]

อัตราส่วนคาบของลูกตุ้ม

\[
\frac{T_2}{T_1}
=
\sqrt{
\frac{\ell_2}{\ell_1}
\frac{g_1}{g_2}
}
\]

การขยายตัวเชิงความร้อน

\[
\ell'=\ell(1+\alpha\Delta T)
\]

\[
\frac{T'}{T}
=
\sqrt{1+\alpha\Delta T}
\]
\end{tcolorbox}

% =========================================================
\section{ข้อสอบจริง สอวน. ที่ใช้ความรู้เรื่องคาบของ SHM}



% =========================================================
\subsection{ปี 2562 ข้อ 8: คาบของลูกตุ้มและการขยายตัวของลวด}

\vspace{1.5em}

\IfFileExists{img/posn62q8.png}{
\begin{center}
    \includegraphics[width=1\linewidth]{img/posn62q8.png}
\end{center}
}{
\begin{tcolorbox}[title=ตำแหน่งภาพที่แนะนำ]
ให้ตัดภาพข้อสอบปี 2562 ข้อ 8 แล้วบันทึกเป็น

\[
\texttt{img/posn62q8.png}
\]
\end{tcolorbox}
}

\begin{examplebox}{เฉลยละเอียด}
คาบของลูกตุ้มอย่างง่ายคือ

\[
T=2\pi\sqrt{\frac{\ell}{g}}
\]

ดังนั้น

\[
T\propto\sqrt{\ell}
\]

เมื่อลวดมีอุณหภูมิสูงขึ้น

\[
\Delta T
\]

และมีสัมประสิทธิ์การขยายตัวเชิงเส้น

\[
\alpha
\]

ความยาวใหม่ของลวดคือ

\[
\ell'=\ell(1+\alpha\Delta T)
\]

อัตราส่วนคาบใหม่ต่อคาบเดิมคือ

\[
\frac{T'}{T}
=
\sqrt{\frac{\ell'}{\ell}}
\]

แทนค่า

\[
\frac{T'}{T}
=
\sqrt{
\frac{\ell(1+\alpha\Delta T)}{\ell}
}
\]

ดังนั้น

\[
\boxed{
\frac{T'}{T}
=
\sqrt{1+\alpha\Delta T}
}
\]

\[
\boxed{\text{ตอบข้อ 3}}
\]
\end{examplebox}

% =========================================================
\subsection{ปี 2562 ข้อ 24: ช่วงเวลาระหว่างการชนของลูกตุ้ม}

\vspace{1.5em}

\IfFileExists{img/posn62q24.png}{
\begin{center}
    \includegraphics[width=1\linewidth]{img/posn62q24.png}
\end{center}
}{
\begin{tcolorbox}[title=ตำแหน่งภาพที่แนะนำ]
ให้ตัดภาพข้อสอบปี 2562 ข้อ 24 แล้วบันทึกเป็น

\[
\texttt{img/posn62q24.png}
\]
\end{tcolorbox}
}

\begin{examplebox}{เฉลยละเอียด}
ลูกตุ้มทั้งสองมีความยาวเท่ากัน

\[
\ell
\]

ดังนั้นลูกตุ้มทั้งสองมีคาบเท่ากัน

\[
T=2\pi\sqrt{\frac{\ell}{g}}
\]

หลังจากชนกันครั้งแรกที่ตำแหน่งต่ำสุด ลูกตุ้มทั้งสองไม่ติดกัน และแยกออกจากตำแหน่งต่ำสุด

สำหรับลูกตุ้มอย่างง่าย ลูกตุ้มแต่ละลูกจะกลับมาถึงตำแหน่งต่ำสุดอีกครั้งหลังเวลาครึ่งคาบ

\[
\Delta t=\frac{T}{2}
\]

แทนสูตรคาบ

\[
\Delta t
=
\frac{1}{2}
\left(
2\pi\sqrt{\frac{\ell}{g}}
\right)
\]

ดังนั้น

\[
\boxed{
\Delta t
=
\pi\sqrt{\frac{\ell}{g}}
}
\]
\end{examplebox}

